
\documentclass{report}

\input{../../preamble}
\input{../../macros}
\input{../../letterfonts}

\title{}
\author{Ryan Lin}
\date{}

\begin{document}
\maketitle
\section*{T/F}

\begin{enumerate}
	\item For the Gilbert Model of Transmission, the steady state distribution can be calculated by $ \pi = (\frac{q}{p+q}, \frac{p}{p+q})  $
	\item The steady state distribution of the doubly stochastic Markov Chain is uniform. 
	\item If $ \{X_1, X_2, \dots, X_{n}\} $ is an irreducible, positive recurrent chain, and $ r $ is bounded or linear, then this is a necessary condition for the existence of a LRAC
	\item An Exponential distribution is memoryless. 
	\item It is possible to exit a recurrent class.
\end{enumerate}

\subsection*{Key}
\begin{enumerate}
	\item True
	\item True
	\item False, it is a sufficient condition
	\item True, use the survival function 
	\item False, you stay in a recurrent class by definition
\end{enumerate}

\section*{Computational}

The Department is considering a proposal to sell IEOR merch and would like to determine the optimal inventory model for backpacks to minimize Long Run Average Costs. They have a storage limit of 3 parcels of backpacks, and would want to know when they should order new backpacks. They are considering either adopting a (0, 3) or (1,3) policy. They project that daily demand $ D_{i} \sim (.2, .5, .3)$ is i.i.d., and that their holding costs is \$4 per parcel per day. They also have a \$10 stockout cost if a customer orders a backpack but they do not have any in stock.  

\begin{enumerate}
	\item What is the expected number of parcels in stock? 
	\item What is the expected demand? 
	\item Which policy would minimize LRAC? 
\end{enumerate}

\subsection*{Key}

\begin{enumerate}
    \item \textbf{Expected Demand:} 
    \[ E[D] = 0(0.2) + 1(0.5) + 2(0.3) = 1.1 \text{ parcels.} \]

    \item \textbf{Expected Parcels in Stock (Start of Day):}
    \begin{itemize}
        \item For $(0,3)$, $\pi = [\frac{25}{77}, \frac{24}{77}, \frac{28}{77}]$, $E[X] \approx 2.04$.
        \item For $(1,3)$, $\pi = [0, \frac{5}{13}, \frac{8}{13}]$, $E[X] \approx 2.62$.
    \end{itemize}

    The LRAC is defined as $E[\text{Holding Costs}] + E[\text{Stockout Costs}]$.
    \begin{itemize}
        \item \textbf{LRAC for (0,3):} 
        Expected Stockout $E[SC] = \pi_1 \cdot 10 \cdot P(D=2) = \frac{25}{77} \cdot 10 \cdot 0.3 \approx 0.97$.
        Total $LRAC = 3.76 + 0.97 = \$4.73$.
        
        \item \textbf{LRAC for (1,3):} 
        Expected Stockout $E[SC] = 0$ (since $X_{start} \ge 2$ and $D_{max} = 2$).
        Total $LRAC = 6.03 + 0 = \$6.03$.
    \end{itemize}
The (0, 3) policy minimizes the Long Run Average Cost.
\end{enumerate}
\begin{align*}
	E[T] = \frac{1}{\lambda} + \frac{\lambda}{\lambda + \mu } \cdot \frac{1}{\lambda} + \frac{\lambda}{\lambda +\mu }E[T]
\end{align*}
\end{document}
